\section{Discussion, Competition Governance, and Repository Lifecycle}

A central operational question addressed in this milestone is the relationship between our Python competition submission notebooks (Track 1) and the standalone \texttt{NeuroARC.jl} reference package (Track 3).

\subsection{Empirical Context and Decontamination Governance}
Three engineering constraints shaped the lifecycle from internal development to open-source release:

\begin{enumerate}
  \item \textbf{Audit and Specification Invariants}: The anti-solipsism framework, the deterministic verifier threat model, and the constructive certificate calculus separate verified numerical components (\texttt{VERIFIED\_NUMERICAL\_COMPONENT}) from unverified empirical hypotheses.
  \item \textbf{Kaggle ARC-AGI-2 Decontamination Governance}: In our audit of the upstream Python \texttt{dsl.py} module (164 top-level \texttt{FunctionDef} nodes), 73 definitions (72 unique runtime keys due to a duplicate \texttt{solve\_b9e38dc0} definition) were task-keyed \texttt{solve\_<8hex>} lookup tables rather than generalizable operators, leaving 91 reusable grid helpers. Exact-match regex stripping in the submitted clean run removed 71 runtime keys and leaked one renamed alias (\texttt{solve\_9385bd28\_generalized}) through a default-argument binding, collapsing accuracy from 39.17 on the public evaluation set (scored locally) to 0.00 on the hidden leaderboard. That failure motivated the constructive verification gates formalized here.
  \item \textbf{Cross-Language Boundary between Python and Julia}: The Kaggle ARC-AGI-2 notebooks run in Python on Kaggle's offline four-L4 GPU environment within the twelve-hour wall-clock limit (reaching 31.94 and 29.72 under our triple-pass NVARC test-time training schedule, and 30.28 on S1, 28.47 on the triangulated replicate, and also 28.47 on S5 with the Python sidecar). \texttt{NeuroARC.jl} is the standalone Julia formalization of the sidecar's verification and typed-composition contract rather than the runtime engine executed inside the Kaggle submission kernels.
\end{enumerate}

\subsection{Open-Source Verification and Reproducibility}
For the ARC Prize 2026 Paper Track release, the repository is packaged as a clean, zero-dependency Julia codebase verified by:
\begin{enumerate}
  \item Deterministic execution of all 652 unit-test assertions via \texttt{Pkg.test()} across core types, constructive proof certificates, preflight grafting contracts, and text-cleanliness invariants.
  \item Automated Graphviz and Tectonic compilation of the manuscript and architectural diagrams (\texttt{Figure B1} and \texttt{Figure B2}).
  \item Public release alongside the companion Kaggle submission notebooks so reviewers can inspect both the empirical Python test-time training schedule and the formal Julia verification contract.
\end{enumerate}
